\documentclass[12pt, a4paper]{article}

% ── Packages ──────────────────────────────────────────────────────────────────
\usepackage[T1]{fontenc}
\usepackage[utf8]{inputenc}
\usepackage[margin=2.6cm, top=3.0cm, bottom=3.0cm]{geometry}
\usepackage{lmodern}
\usepackage{microtype}
\usepackage{parskip}
\usepackage{abstract}
\usepackage{booktabs}
\usepackage{array}
\usepackage{hyperref}
\usepackage{tikz}
\usetikzlibrary{arrows.meta, positioning, shapes.rounded}

\hypersetup{colorlinks=true, urlcolor=black, linkcolor=black,
            pdftitle={dedomena -- Internal Product Overview}}

% ── Title block ───────────────────────────────────────────────────────────────
\title{%
  \textbf{dedomena}\\[0.4em]
  \large Internal Product Overview
}
\author{%
  Made by Sami El-Figha\\[0.2em]
  \small\href{https://linkedin.com/in/samielgha}{linkedin.com/in/samielgha}
}
\date{April 7, 2026}

% ── Document ──────────────────────────────────────────────────────────────────
\begin{document}

\maketitle
\thispagestyle{plain}

% ── Abstract ──────────────────────────────────────────────────────────────────
\begin{abstract}
In most organisations, data questions still pass through a small team of analysts,
which means waiting hours or even days for answers that should be immediate.
dedomena solves this by letting any team member type a question in plain English
and receive a structured, sourced answer in seconds, with no knowledge of SQL or
data architecture required. In the background, a three-tier pipeline handles
everything: a question parser determines the real intent of the query, a data
retriever extracts only the relevant portion from connected sources (CRM,
databases, spreadsheets, finance tools, and more), and an answer builder formats
the result and flags any detected anomalies. Raw data never leaves the
organisation's infrastructure, every query is logged for audit purposes, and
analysts get their time back for work that actually requires their expertise.
\end{abstract}

% ── System Architecture diagram ───────────────────────────────────────────────
\section*{System Architecture}

\begin{figure}[h]
\centering
\begin{tikzpicture}[
    node distance=0.9cm and 1.4cm,
    tier/.style={font=\footnotesize\itshape, text=black!50},
    box/.style={rectangle, rounded corners=4pt, draw=black!40,
                fill=black!4, font=\small, inner sep=8pt,
                minimum height=0.9cm, minimum width=2.8cm, text centered},
    src/.style={rectangle, rounded corners=4pt, draw=black!30,
                fill=black!2, font=\small, inner sep=7pt,
                minimum height=0.8cm, minimum width=2.6cm, text centered},
    arr/.style={-{Stealth[length=5pt]}, thick, black!50}
  ]

  % Tier labels (left margin)
  \node[tier] (t1) at (-1.0, 0)    {Tier 1};
  \node[tier] (t2) at (-1.0, -1.9) {Tier 2};
  \node[tier] (t3) at (-1.0, -4.0) {Tier 3};

  % Tier 1 -- User query
  \node[box] (query) at (4.5, 0) {Query};

  % Tier 2 -- Processing pipeline
  \node[box] (parser)  at (1.5, -1.9) {Parser};
  \node[box] (ret)     at (4.5, -1.9) {Retriever};
  \node[box] (builder) at (7.5, -1.9) {Answer Builder};

  % Tier 3 -- Data sources
  \node[src] (crm)   at (1.5, -4.0) {CRM / Sales};
  \node[src] (db)    at (4.5, -4.0) {Databases};
  \node[src] (files) at (7.5, -4.0) {Files / Finance};

  % Arrows: Query down to pipeline
  \draw[arr] (query) -- (parser);
  \draw[arr] (query) -- (ret);
  \draw[arr] (query) -- (builder);

  % Arrows: pipeline left to right
  \draw[arr] (parser)  -- (ret);
  \draw[arr] (ret)     -- (builder);

  % Arrows: sources up to retriever
  \draw[arr] (crm)   -- (ret);
  \draw[arr] (db)    -- (ret);
  \draw[arr] (files) -- (ret);

\end{tikzpicture}
\caption{Three-tier system architecture.}
\end{figure}

% ── The Problem ───────────────────────────────────────────────────────────────
\section*{The Problem}

Getting a data answer today takes one to three business days. A manager emails an
analyst, the analyst locates data across multiple systems, writes a query, formats
a report, and sends it back. By then the decision context has often changed.

The root cause is fragmentation. A question like \textit{Which region grew fastest
this quarter?} may require pulling from a CRM, a finance system, and several
spreadsheets at once. No conventional tool bridges those sources. Every cross-system
question becomes manual work, and a small number of technical staff become a
bottleneck for the whole organisation.

% ── The Solution ──────────────────────────────────────────────────────────────
\section*{The Solution}

dedomena connects to all data sources once and lets any authorised user type a
question in plain English. The platform queries all relevant sources simultaneously
and returns a structured, sourced answer in under 30 seconds. No SQL. No formulas.
No waiting.

The workflow reduces to three steps:

\begin{enumerate}
  \item \textbf{Connect.} An administrator links data sources once. No ongoing
        maintenance required.
  \item \textbf{Ask.} Any user types a question. The platform handles source
        selection and retrieval automatically.
  \item \textbf{Act.} A sourced answer is returned in the most useful format:
        prose, table, or list.
\end{enumerate}

% ── How It Works ──────────────────────────────────────────────────────────────
\section*{How It Works}

\textbf{Question Parser.}
Extracts intent, entities, time ranges, and metrics from the user's input
regardless of phrasing.

\textbf{Data Retriever.}
Queries only the sources relevant to that intent. Raw data never leaves the
organisation's infrastructure. Only the query-relevant excerpt moves forward.

\textbf{Answer Builder.}
Synthesises retrieved excerpts into a response in the most useful format and
flags anomalies automatically.

% ── Supported Data Sources ────────────────────────────────────────────────────
\section*{Supported Data Sources}

\begin{table}[h]
\centering
\begin{tabular}{>{\bfseries}ll}
  \toprule
  Category & Examples \\
  \midrule
  CRM and sales          & Salesforce, HubSpot, Pipedrive \\
  Databases              & PostgreSQL, MySQL, BigQuery, Snowflake \\
  Files and spreadsheets & Excel, CSV, PDF, Google Sheets \\
  Cloud storage          & Google Drive, OneDrive, SharePoint \\
  Finance                & QuickBooks, Stripe \\
  Project management     & Notion, Airtable, Jira, Asana \\
  \bottomrule
\end{tabular}
\caption{Data source categories supported at initial release. Over 100 source
types in total.}
\end{table}

% ── Who Benefits ──────────────────────────────────────────────────────────────
\section*{Who Benefits}

\begin{table}[h]
\centering
\begin{tabular}{>{\bfseries}lp{9cm}}
  \toprule
  Role & What they get \\
  \midrule
  Management  & Answers to performance questions without waiting for reports. \\[4pt]
  Sales       & Account history and pipeline health on demand. \\[4pt]
  Operations  & A unified view across inventory, suppliers, and process data. \\[4pt]
  Finance     & Cross-system budget and variance checks without formulas. \\[4pt]
  Analysts    & Freed from routine queries. Focus on higher-order analysis. \\
  \bottomrule
\end{tabular}
\caption{Roles and corresponding benefits.}
\end{table}

% ── Data Privacy ──────────────────────────────────────────────────────────────
\section*{Data Privacy}

Raw data never leaves the organisation's environment. When a question is asked,
only the query-relevant excerpt is passed to the reasoning layer. Everything else
stays in place. This satisfies data residency requirements in regulated industries.

% ── Expected Benefits ─────────────────────────────────────────────────────────
\section*{Expected Benefits}

\begin{enumerate}
  \item \textbf{Speed.} One to three day turnaround collapses to under 30 seconds.
  \item \textbf{Accessibility.} No technical training required. Any staff member
        can use it from day one.
  \item \textbf{Single interface.} One place to query all sources instead of
        logging into multiple systems.
  \item \textbf{Audit trail.} Every query is logged with the sources consulted
        and a timestamp.
  \item \textbf{Analyst capacity.} Technical staff are freed from routine queries
        and redirected to deeper work.
\end{enumerate}

% ── Rollout Plan ──────────────────────────────────────────────────────────────
\section*{Rollout Plan}

\begin{table}[h]
\centering
\begin{tabular}{>{\bfseries}lll}
  \toprule
  Phase & When & Activities \\
  \midrule
  1\ \ Pilot    & Weeks 1--2  & Connect 2--3 key sources. Trial with managers. Collect feedback. \\[4pt]
  2\ \ Rollout  & Month 2     & Connect all sources. Open to all teams. Short onboarding session. \\[4pt]
  3\ \ Optimise & Month 3+    & Review common queries. Build standing summaries. \\
  \bottomrule
\end{tabular}
\caption{Recommended phased rollout plan.}
\end{table}

% ── Conclusion ────────────────────────────────────────────────────────────────
\section*{Conclusion}

dedomena gives any member of the organisation the ability to ask a question about
business data and get a clear, sourced answer in seconds, without writing code,
without waiting for a report, and without knowing where the data lives. It removes
the bottleneck on technical staff and gives every team direct access to the
information they need.

\vspace{1.5em}
\noindent\textcolor{black!40}{\small This document is intended for internal review.}

\end{document}
